\section{Canonical Organisation and Workspace Authority for AION Flow}
\label{sec:aion-flow-canonical-authority}

This section records completion of \texttt{AFLOW-1104}. It is intended to be appended to the
existing AION and Pilot technical document. AION Flow no longer accepts a role name supplied by
the browser as evidence that a person may inspect, author, export or commit a workflow.

\subsection{Canonical decision path}

Every protected production operation now follows this decision sequence:

\begin{quote}
Signed-in person $\rightarrow$ active workspace lookup $\rightarrow$ canonical organisation
membership $\rightarrow$ exact capability decision $\rightarrow$ workflow scope check
$\rightarrow$ operation or fail-closed refusal.
\end{quote}

The workspace repository proves that the requested workspace exists and is active. The existing
Organisation Authority container then proves that the person is active and that one of their
persisted roles grants the exact AION Flow capability. A browser may still include a role label for
display or old-client compatibility, but the backend marks it as ignored and never uses it to
grant access.

The initial exact capabilities cover review, authoring, signed commits, encrypted export,
encrypted import and benchmarking. They use the \texttt{aion\_flow} capability namespace and are
included in the canonical Owner/Director template. An
organisation administrator may later assign individual capabilities to custom roles through the
same authority model; no separate AION Flow permissions database is created.

\subsection{Protected surfaces}

Canonical authority is enforced for production qualification, signed revision commits, revision
history, version differences, encrypted export, encrypted import, zero-content telemetry and local
benchmarks. The desktop Production panel sends the canonical person and workspace identity. It no
longer overwrites that identity with an \texttt{owner} role claim. The qualification result displays
the bounded authority decision so the user can see which store made the decision and why access was
allowed or denied.

The result fails closed when the person identifier is absent, the workspace does not exist or is
inactive, the person is not active, the capability is not granted, or the requested organisation,
workspace or graph scope does not match. Department-scoped graphs also pass their department to the
existing authority evaluator.

\subsection{Tenant isolation and revision ownership}

Workflow revision storage is keyed by both workspace and workflow identifier. Two businesses may
therefore use the same workflow name without sharing history or colliding revisions. Each new signed
revision records the canonical workspace, canonical organisation, person identifier and authority
source alongside its graph hash, optimistic base revision, timestamp, public key and signature.
The person's mutable browser role is deliberately absent from the signed authorship record.

Historical revision ledgers created before workspace scoping are not guessed or silently assigned.
They fail closed with an explicit migration requirement. This prevents a member of one workspace
from claiming an ambiguous legacy history merely by knowing its workflow identifier.

\subsection{Verification and remaining work}

Focused automated tests prove that a forged \texttt{owner} claim cannot elevate an ordinary
employee, a missing workspace is rejected, canonical owners retain their intended capabilities,
the same workflow identifier remains isolated across two active workspaces, signed revision
conflicts still fail, and encrypted transfer, telemetry and benchmarking remain governed.

This completes the membership and workspace-authority integration. It does not complete overall
enterprise qualification. Formal threat and dependency review, sustained parallel and long-running
benchmarks, accessibility and regulatory review, and representative business field testing remain
open under \texttt{AFLOW-1107}--\texttt{AFLOW-1109}.
